%! Interpreting p-Values — Unit 6, Lesson 6.5 — Exit Ticket (Answer Key)
\documentclass[10pt]{article}
\usepackage{apstats-article}
\usepackage{apstats-key}   % pulls in -boxes; do NOT also load -boxes
\newcommand{\TallMath}[1]{$\displaystyle #1\rule[-1.4em]{0pt}{3.2em}$}

\begin{document}

\pageheader{Unit 6, Lesson 6.5}{Exit Ticket --- Answer Key}
\namedateperiod

\vspace{0.15in}

A researcher claims that 25\% of adults report daily physical activity. $z \approx 1.20$.

\begin{enumerate}
  \item $H_0$: \ans{$p = 0.25$} \quad $H_a$: \ans{$p > 0.25$}

  \item Tail: \ans{right} \quad $p$-value $= P(Z \ge 1.20) = $ \ans{$1 - 0.8849 = 0.1151$}

  \item \ansline{Assuming 25\% of adults report daily physical activity, there is an 11.5\% probability of observing a sample proportion of $84/300=0.28$ or higher by chance alone.}

  \item \ansline{Error: the $p$-value is not the probability $H_0$ is true. Correction: the $p$-value is the probability of getting a result this extreme \emph{if} $H_0$ were true, not the probability that $H_0$ actually is true.}
\end{enumerate}

\end{document}
